% ============================================================
% EJERCICIO 1 — Ecuaciones de electrónica (Módulos 1 y 2)
% Objetivo : escribir 6 ecuaciones con buenas prácticas de LaTeX.
% Tareas   : 1) Genera el código con una IA (guarda el prompt abajo).
%            2) Complete cada TODO. 3) Compile sin errores ni avisos.
% Entregable: este archivo + PDF.      Rúbrica (20 pts) al final.
% PROMPT UTILIZADO (obligatorio): TODO pega aquí tu prompt
% ============================================================
\documentclass[12pt,a4paper]{article}
\usepackage[utf8]{inputenc}
\usepackage[T1]{fontenc}
\usepackage[spanish]{babel}
\usepackage{amsmath,amssymb}
\usepackage{booktabs}
\usepackage{siunitx}
\sisetup{locale=FR, output-decimal-marker={,}}
\usepackage{circuitikz}
\usepackage[hidelinks]{hyperref}

\title{Ecuaciones de electrónica}
\author{TODO Nombre y apellidos}
\date{\today}

\begin{document}
\maketitle

\section{Ecuaciones fundamentales}
\begin{enumerate}
  \item Ley de Ohm: % TODO  V = I \cdot R, con \begin{equation} y \label
  \item Potencia eléctrica (tres formas equivalentes): % TODO
  \item Reactancia capacitiva $X_C$: % TODO fracción con \frac
  \item Reactancia inductiva $X_L$: % TODO
  \item Impedancia RLC serie: % TODO raíz con \sqrt
  \item Frecuencia de resonancia $f_0$: % TODO
\end{enumerate}

\section{Cálculo}
% TODO Con $R=\SI{100}{\ohm}$, $L=\SI{10}{\milli\henry}$ y $C=\SI{1}{\micro\farad}$
% calcule $f_0$ y escríbalo con siunitx. Cite una ecuación con \eqref{...}.

\end{document}
% ------------------------------------------------------------
% RÚBRICA (20 pts): compila sin errores (4) · 6 ecuaciones correctas (6)
% · equation + label + \eqref (3) · siunitx en el cálculo (3)
% · prompt documentado y revisado críticamente (4)
% ------------------------------------------------------------
